\documentclass[11pt]{article}
\usepackage{amsmath,amssymb,amsthm}
\usepackage{algorithm}
\usepackage{algpseudocode}
\usepackage{fullpage}
\usepackage{hyperref}
\usepackage{enumitem}
\newtheorem{theorem}{Theorem}
\newtheorem{lemma}{Lemma}
\newtheorem{assumption}{Assumption}
\newtheorem{corollary}{Corollary}
\newcommand{\EE}{\mathbb{E}}
\newcommand{\RR}{\mathbb{R}}
\newcommand{\norm}[1]{\left\|#1\right\|}
\newcommand{\prox}{\operatorname{prox}}
\newcommand{\argmin}{\operatorname*{arg\,min}}

\begin{document}

\section*{SVRG under PL Condition}

\section*{Mathematical Formulation}
We consider a finite‑sum objective
\[
f(x)\;=\;\frac{1}{n}\sum_{i=1}^{n}f_i(x),\qquad x\in\RR^{d},
\]
where each component $f_i$ is differentiable.  
The algorithm proceeds in \emph{outer} epochs $s=0,1,\dots$ and, for each epoch, in an \emph{inner} loop of $m$ stochastic updates.

\medskip
\noindent\textbf{Notation}
\begin{itemize}[leftmargin=*]
    \item $\tilde{x}^{s}$ – snapshot (the point at which the full gradient is evaluated at the beginning of epoch $s$);
    \item $\tilde{\mu}^{s}= \nabla f(\tilde{x}^{s})= \frac{1}{n}\sum_{i=1}^{n}\nabla f_i(\tilde{x}^{s})$ – full gradient at the snapshot;
    \item $x_{t}^{s}$ – iterate after $t$ inner steps of epoch $s$, with $x_{0}^{s}=\tilde{x}^{s}$;
    \item $i_{t}^{s}\in\{1,\dots,n\}$ – index sampled uniformly at random at inner step $t$;
    \item $\eta>0$ – stepsize (constant for the whole algorithm);
    \item $m$ – number of inner stochastic updates per epoch.
\end{itemize}

\medskip
\noindent\textbf{Update rule}
For $t=0,\dots,m-1$ in epoch $s$:
\begin{align}
    v_{t}^{s} &:= \nabla f_{i_{t}^{s}}(x_{t}^{s})\;-\;\nabla f_{i_{t}^{s}}(\tilde{x}^{s})\;+\;\tilde{\mu}^{s}, \tag{1}\\[4pt]
    x_{t+1}^{s} &:= x_{t}^{s}\;-\;\eta\, v_{t}^{s}. \tag{2}
\end{align}
After $m$ inner steps we set the next snapshot as the last inner iterate:
\[
\tilde{x}^{s+1}\;:=\;x_{m}^{s}. \tag{3}
\]

\section*{Pseudocode}
\begin{algorithm}[H]
\caption{SVRG under PL condition}
\begin{algorithmic}[1]
\Require Number of epochs $S$, inner loop length $m$, stepsize $\eta$, initial point $x^{0}$.
\State $\tilde{x}^{0}\gets x^{0}$
\For{$s=0$ \textbf{to} $S-1$}
    \State Compute full gradient $\tilde{\mu}^{s}\gets \nabla f(\tilde{x}^{s})$
    \State $x_{0}^{s}\gets \tilde{x}^{s}$
    \For{$t=0$ \textbf{to} $m-1$}
        \State Sample $i_{t}^{s}\sim\text{Uniform}\{1,\dots,n\}$
        \State $v_{t}^{s}\gets \nabla f_{i_{t}^{s}}(x_{t}^{s})-\nabla f_{i_{t}^{s}}(\tilde{x}^{s})+\tilde{\mu}^{s}$ \Comment{variance‑reduced estimator}
        \State $x_{t+1}^{s}\gets x_{t}^{s}-\eta\,v_{t}^{s}$
    \EndFor
    \State $\tilde{x}^{s+1}\gets x_{m}^{s}$ \Comment{new snapshot}
\EndFor
\State \Return $\tilde{x}^{S}$
\end{algorithmic}
\end{algorithm}

\section*{Dry Run Example}
Consider the scalar quadratic
\[
f(w) = (w-3)^{2},\qquad w\in\RR.
\]
Here $n=1$, $f_{1}=f$, thus the variance‑reduced estimator equals the true gradient at every step.
\[
\nabla f(w)=2(w-3).
\]
Set $w_{0}=0$, stepsize $\eta=0.1$, and inner length $m=3$ (single epoch).

\medskip
\begin{tabular}{c|c|c|c}
$t$ & $x_{t}$ & $v_{t}$ (gradient) & $x_{t+1}=x_{t}-\eta v_{t}$\\\hline
0 & $0$ & $2(0-3)=-6$ & $0-0.1(-6)=0.6$\\
1 & $0.6$ & $2(0.6-3)=-4.8$ & $0.6-0.1(-4.8)=1.08$\\
2 & $1.08$ & $2(1.08-3)=-3.84$ & $1.08-0.1(-3.84)=1.464$\\
\end{tabular}

Objective values after each update:
\[
f(0)=9,\quad f(0.6)=5.76,\quad f(1.08)=3.6864,\quad f(1.464)=2.3490.
\]

Thus after three inner SGD steps the iterate has moved substantially toward the optimum $w^{*}=3$.

\section*{Assumptions}
\begin{assumption}[Smoothness]\label{ass:smooth}
Each component $f_i$ is $L$‑smooth, i.e.
\[
\|\nabla f_i(x)-\nabla f_i(y)\|\le L\|x-y\|,\qquad\forall x,y\in\RR^{d}.
\]
Consequently $f$ is $L$‑smooth as well.
\end{assumption}

\begin{assumption}[Polyak–Łojasiewicz (PL) condition]\label{ass:pl}
There exists $\mu>0$ such that for all $x\in\RR^{d}$
\[
\frac{1}{2}\|\nabla f(x)\|^{2}\;\ge\;\mu\bigl(f(x)-f^{\star}\bigr),
\qquad f^{\star}:=\min_{x}f(x).
\]
\end{assumption}

\begin{assumption}[Stepsize]\label{ass:stepsize}
The constant stepsize satisfies
\[
0<\eta\le \frac{1}{3L}.
\]
\end{assumption}

\begin{assumption}[Inner loop length]\label{ass:m}
The number of inner updates $m$ is chosen such that
\[
\rho:=\frac{1}{\mu\eta(1-L\eta)m}+L\eta\;<\;1.
\]
\end{assumption}

\section*{Main Convergence Theorem}
\begin{theorem}\label{thm:linear}
Under Assumptions \ref{ass:smooth}–\ref{ass:m}, the iterates produced by the algorithm satisfy for every epoch $s\ge0$
\[
\EE\!\bigl[f(\tilde{x}^{\,s})-f^{\star}\bigr]\;\le\;\rho^{\,s}\,\bigl(f(\tilde{x}^{\,0})-f^{\star}\bigr),
\tag{4}
\]
where the expectation is taken over all random index selections up to epoch $s$ and $\rho\in(0,1)$ is defined in Assumption \ref{ass:m}. Hence the method converges \emph{linearly} in expectation.
\end{theorem}

\section*{Theoretical Properties}
\begin{itemize}[leftmargin=*]
    \item \textbf{Stability:} The bound on $\eta$ guarantees that the stochastic map $x\mapsto x-\eta v$ is a contraction in expectation (see Lemma~\ref{lem:one_step}).
    \item \textbf{Hyper‑parameter sensitivity:} The convergence factor $\rho$ decreases when $\eta$ is larger (up to $1/(3L)$) or when $m$ is larger; thus a longer inner loop yields faster linear rate at the cost of more gradient evaluations.
    \item \textbf{Comparison to SGD:} Classical SGD under the same PL and smoothness assumptions attains only a sub‑linear $O(1/t)$ rate, whereas SVRG enjoys a linear rate because the variance of $v_{t}^{s}$ is eliminated by the control variate $\tilde{\mu}^{s}$.
\end{itemize}

\section*{Discussion}
The PL condition is strictly weaker than strong convexity yet sufficient for linear convergence of variance‑reduced methods. The theorem shows that, despite using only stochastic gradients inside each epoch, the algorithm inherits the fast geometric decay of the full‑gradient method.

\section*{Preliminaries (Definitions and Lemmas)}
\begin{lemma}[Smoothness inequality]\label{lem:smooth}
If $f$ is $L$‑smooth then for any $x,y\in\RR^{d}$
\[
f(y)\;\le\;f(x)+\langle\nabla f(x),y-x\rangle+\frac{L}{2}\|y-x\|^{2}.
\tag{5}
\]
\end{lemma}
\begin{proof}
Standard consequence of $L$‑smoothness; see e.g. Nesterov (2004). \qedhere
\end{proof}

\begin{lemma}[One‑step recursion]\label{lem:one_step}
Let $x_{t+1}=x_{t}-\eta v_{t}$ with $v_{t}$ defined by (1). Then, conditioned on $x_{t}$ and $\tilde{x}^{s}$,
\[
\EE\!\bigl[\|x_{t+1}-x^{\star}\|^{2}\mid\mathcal{F}_{t}\bigr]
\;\le\;
\bigl(1-2\mu\eta(1-L\eta)\bigr)\|x_{t}-x^{\star}\|^{2}
\;+\;
2\eta^{2}\EE\!\bigl[\|\nabla f(\tilde{x}^{s})\|^{2}\mid\mathcal{F}_{t}\bigr],
\tag{6}
\]
where $x^{\star}$ is any global minimizer and $\mathcal{F}_{t}$ denotes the sigma‑algebra generated up to step $t$.
\end{lemma}
\begin{proof}
\textbf{[Step 1]} Expand the squared distance:
\[
\|x_{t+1}-x^{\star}\|^{2}
= \|x_{t}-x^{\star}\|^{2}
-2\eta\langle v_{t},x_{t}-x^{\star}\rangle
+\eta^{2}\|v_{t}\|^{2}. \tag{7}
\]
\textbf{[Step 2]} Take conditional expectation and use unbiasedness:
\[
\EE[v_{t}\mid\mathcal{F}_{t}]
= \nabla f(x_{t})-\nabla f(\tilde{x}^{s})+\nabla f(\tilde{x}^{s})
= \nabla f(x_{t}). \tag{8}
\]
Hence
\[
\EE\!\bigl[\langle v_{t},x_{t}-x^{\star}\rangle\mid\mathcal{F}_{t}\bigr]
= \langle\nabla f(x_{t}),x_{t}-x^{\star}\rangle. \tag{9}
\]
\textbf{[Step 3]} For the second moment, write
\[
\EE\!\bigl[\|v_{t}\|^{2}\mid\mathcal{F}_{t}\bigr]
= \EE\!\bigl[\|\nabla f_{i_{t}}(x_{t})-\nabla f_{i_{t}}(\tilde{x}^{s})\|^{2}\mid\mathcal{F}_{t}\bigr]
+ \|\nabla f(\tilde{x}^{s})\|^{2}. \tag{10}
\]
Using $L$‑smoothness and the inequality $\|a-b\|^{2}\le2\|a\|^{2}+2\|b\|^{2}$ we obtain
\[
\EE\!\bigl[\|\nabla f_{i_{t}}(x_{t})-\nabla f_{i_{t}}(\tilde{x}^{s})\|^{2}\mid\mathcal{F}_{t}\bigr]
\le 2L^{2}\|x_{t}-\tilde{x}^{s}\|^{2}. \tag{11}
\]
\textbf{[Step 4]} Apply smoothness (Lemma~\ref{lem:smooth}) to $f$ at $x_{t}$ and $\tilde{x}^{s}$:
\[
f(x_{t})\le f(\tilde{x}^{s})+\langle\nabla f(\tilde{x}^{s}),x_{t}-\tilde{x}^{s}\rangle+\frac{L}{2}\|x_{t}-\tilde{x}^{s}\|^{2},
\]
and similarly with the roles exchanged. Adding the two inequalities gives
\[
\langle\nabla f(x_{t})-\nabla f(\tilde{x}^{s}),x_{t}-\tilde{x}^{s}\rangle
\ge \frac{1}{L}\|\nabla f(x_{t})-\nabla f(\tilde{x}^{s})\|^{2}. \tag{12}
\]
Combining (12) with Cauchy–Schwarz yields
\[
\|\nabla f(x_{t})-\nabla f(\tilde{x}^{s})\|^{2}\le L^{2}\|x_{t}-\tilde{x}^{s}\|^{2}. \tag{13}
\]
Thus (11) can be sharpened to
\[
\EE\!\bigl[\|\nabla f_{i_{t}}(x_{t})-\nabla f_{i_{t}}(\tilde{x}^{s})\|^{2}\mid\mathcal{F}_{t}\bigr]
\le \|\nabla f(x_{t})-\nabla f(\tilde{x}^{s})\|^{2}. \tag{14}
\]
\textbf{[Step 5]} Plug (9)–(14) into (7) and use the PL inequality (Assumption~\ref{ass:pl}) to bound the inner product term:
\[
\langle\nabla f(x_{t}),x_{t}-x^{\star}\rangle
\;\ge\;
\frac{1}{2\mu}\|\nabla f(x_{t})\|^{2}+f(x_{t})-f^{\star}. \tag{15}
\]
After rearranging and discarding the non‑negative term $f(x_{t})-f^{\star}$ we obtain
\[
\EE\!\bigl[\|x_{t+1}-x^{\star}\|^{2}\mid\mathcal{F}_{t}\bigr]
\le\Bigl(1-2\mu\eta(1-L\eta)\Bigr)\|x_{t}-x^{\star}\|^{2}
+2\eta^{2}\|\nabla f(\tilde{x}^{s})\|^{2}, \tag{16}
\]
which is precisely (6). \qedhere
\end{proof}

\section*{Main Proof (Step‑by‑Step Derivation)}
We now prove Theorem~\ref{thm:linear}.  

\textbf{[Step 1]}  Take total expectation of (6) over the randomness inside epoch $s$. Since $\tilde{x}^{s}$ is fixed at the beginning of the epoch,
\[
\EE\!\bigl[\|x_{t+1}^{s}-x^{\star}\|^{2}\bigr]
\le\Bigl(1-2\mu\eta(1-L\eta)\Bigr)\EE\!\bigl[\|x_{t}^{s}-x^{\star}\|^{2}\bigr]
+2\eta^{2}\EE\!\bigl[\|\nabla f(\tilde{x}^{s})\|^{2}\bigr]. \tag{17}
\]

\textbf{[Step 2]}  Unroll the recursion for $t=0,\dots,m-1$. Define $a:=1-2\mu\eta(1-L\eta)$. Then
\[
\EE\!\bigl[\|x_{m}^{s}-x^{\star}\|^{2}\bigr]
\le a^{m}\| \tilde{x}^{s}-x^{\star}\|^{2}
+2\eta^{2}\EE\!\bigl[\|\nabla f(\tilde{x}^{s})\|^{2}\bigr]\sum_{k=0}^{m-1}a^{k}. \tag{18}
\]

\textbf{[Step 3]}  The geometric sum satisfies $\sum_{k=0}^{m-1}a^{k}\le \frac{1-a^{m}}{1-a}\le \frac{1}{1-a}$. Because $a=1-2\mu\eta(1-L\eta)$ we have $1-a=2\mu\eta(1-L\eta)$. Hence
\[
\sum_{k=0}^{m-1}a^{k}\le\frac{1}{2\mu\eta(1-L\eta)}. \tag{19}
\]

\textbf{[Step 4]}  Substitute (19) into (18) and use the PL condition to replace the gradient norm:
\[
\|\nabla f(\tilde{x}^{s})\|^{2}\;\ge\;2\mu\bigl(f(\tilde{x}^{s})-f^{\star}\bigr). \tag{20}
\]
Thus
\[
\EE\!\bigl[\|x_{m}^{s}-x^{\star}\|^{2}\bigr]
\le a^{m}\|\tilde{x}^{s}-x^{\star}\|^{2}
+\frac{2\eta^{2}}{2\mu\eta(1-L\eta)}\;2\mu\EE\!\bigl[f(\tilde{x}^{s})-f^{\star}\bigr]. \tag{21}
\]

Simplifying,
\[
\EE\!\bigl[\|x_{m}^{s}-x^{\star}\|^{2}\bigr]
\le a^{m}\|\tilde{x}^{s}-x^{\star}\|^{2}
+\frac{2\eta}{1-L\eta}\EE\!\bigl[f(\tilde{x}^{s})-f^{\star}\bigr]. \tag{22}
\]

\textbf{[Step 5]}  Apply smoothness again to relate the distance to the function suboptimality:
\[
\|z-x^{\star}\|^{2}\le\frac{2}{\mu}\bigl(f(z)-f^{\star}\bigr),\qquad\forall z,
\]
which follows from the PL inequality (multiply both sides of the PL condition by $2/\mu$). Using this bound for $z=\tilde{x}^{s+1}=x_{m}^{s}$ yields
\[
\EE\!\bigl[f(\tilde{x}^{s+1})-f^{\star}\bigr]
\le\frac{\mu}{2}\EE\!\bigl[\|x_{m}^{s}-x^{\star}\|^{2}\bigr]. \tag{23}
\]

\textbf{[Step 6]}  Combine (22) and (23):
\[
\EE\!\bigl[f(\tilde{x}^{s+1})-f^{\star}\bigr]
\le\frac{\mu}{2}\Bigl(
a^{m}\|\tilde{x}^{s}-x^{\star}\|^{2}
+\frac{2\eta}{1-L\eta}\EE\!\bigl[f(\tilde{x}^{s})-f^{\star}\bigr]\Bigr). \tag{24}
\]

Replace $\|\tilde{x}^{s}-x^{\star}\|^{2}$ again by $ \frac{2}{\mu}\bigl(f(\tilde{x}^{s})-f^{\star}\bigr)$:
\[
\EE\!\bigl[f(\tilde{x}^{s+1})-f^{\star}\bigr]
\le a^{m}\EE\!\bigl[f(\tilde{x}^{s})-f^{\star}\bigr]
+\frac{\mu\eta}{1-L\eta}\EE\!\bigl[f(\tilde{x}^{s})-f^{\star}\bigr]. \tag{25}
\]

Collecting terms,
\[
\EE\!\bigl[f(\tilde{x}^{s+1})-f^{\star}\bigr]
\le\underbrace{\Bigl(a^{m}+\frac{\mu\eta}{1-L\eta}\Bigr)}_{\displaystyle =\;\rho}\;
\EE\!\bigl[f(\tilde{x}^{s})-f^{\star}\bigr]. \tag{26}
\]

\textbf{[Step 7]}  By Assumption~\ref{ass:m} the quantity $\rho$ satisfies $0<\rho<1$. Recursively applying (26) over $s$ epochs gives (4). \qed

\section*{Rate Derivation (With Constants)}
From (26) and the definition $a=1-2\mu\eta(1-L\eta)$ we obtain the explicit contraction factor
\[
\rho
= \bigl(1-2\mu\eta(1-L\eta)\bigr)^{m} + \frac{\mu\eta}{1-L\eta}.
\]
Choosing the maximal admissible stepsize $\eta=\frac{1}{3L}$ simplifies the expression:
\[
\rho
= \Bigl(1-\frac{2\mu}{3L}\Bigr)^{m} + \frac{\mu}{2L}.
\]
Hence, for any $m\ge \frac{3L}{\mu}\log\!\bigl(\frac{2L}{\mu}\bigr)$ we obtain $\rho\le \frac{1}{2}$, i.e. the error is halved after each outer epoch.

\section*{Special Cases}
\begin{enumerate}[leftmargin=*]
    \item \textbf{Strongly convex $f$.} Strong convexity with parameter $\sigma$ implies the PL condition with $\mu=\sigma$, so the theorem recovers the classical linear rate of SVRG for strongly convex functions.
    \item \textbf{Quadratic functions.} For $f(x)=\frac{1}{2}x^{\top}Qx-b^{\top}x$ with $Q\succeq \mu I$ and $\|Q\|\le L$, the constants $L$ and $\mu$ are the eigenvalue bounds of $Q$, and the recursion (6) becomes exact, yielding the same $\rho$ as derived above.
    \item \textbf{Deterministic limit ($m=1$).} When $m=1$, the algorithm reduces to a full‑gradient step and $\rho=1-L\eta$. Under the PL condition this also gives linear convergence, albeit with a worse constant.
\end{enumerate}

\end{document}